\section{The endpoint gain and the closed surface}
\label{sec:final-proof}

The nonlinear construction used the weak weight $H^\beta$ because
the reference inverse is available for exponents below one. The
metric controls nevertheless acquire twice that exponent. This is
forced by two angular average equations. The first recovers the mean
height from a quadratic expression; the second recovers the metric
controls from that mean and another quadratic expression. There is
no new nonlinear iteration in a stronger function space.
Every height in $\mathcal X_\beta$ already has a smooth flat endpoint
extension. The gain below strengthens the annular result by
quantifying the corrections to the mean height and the metric;
geometric closure uses the flat extensions established in
Section~\ref{sec:fixed-geometry}.

\subsection{The two identities that improve decay}

Let $\mathfrak F_w(U)$ be the connection period from
\eqref{eq:cart-mean-path-flux}. Its derivative at $U=0$ is
$w\av{U_r}$. Write its bilinear Taylor remainder as
\begin{equation}\label{eq:posterior-period-bilinear}
 \mathcal B_{w,U}(v_1,v_2)
 =\int_0^1(1-t)D_U^2\mathfrak F_w(tU)[v_1,v_2]\,dt.
\end{equation}
The normalized period of a solution vanishes, because its radial
derivative is the mean scalar residual and its endpoint value is
zero. Therefore, with $u^{\rm av}=\av u$,
\begin{equation}\label{eq:mean-recovered-derivative}
 (u^{\rm av})'=-q w^{-1}\mathcal B_{w,qu}(u,u).
\end{equation}

For the mean circle map use additive coordinates
$\mathbf v=(Z,\kappa,a)$, where $Z=qu$ and $h_L=\id+a$, and put
\[
 \mathcal C(\mathbf v)=\av{\mathfrak e(\mathbf v)},
 \qquad
 \mathcal B^\circ_{\mathbf v}(v_1,v_2)
 =\int_0^1(1-t)D^2\mathcal C(t\mathbf v)[v_1,v_2]\,dt.
\]
Its first derivative is the identity on $\kappa$, zero on odd chart
directions, and $R_0^\circ$ on the radial mean height. Thus the
second identity at a solution is
\begin{equation}\label{eq:weak-posterior-control}
 0=\kappa+qR_0^\circ u^{\rm av}
                       +\mathcal B^\circ_{\mathbf v}(\mathbf v,\mathbf v).
\end{equation}
The two displayed formulas already explain the exponent $2\beta$.
We must verify these estimates in the stronger topology, including
parameter derivatives, to retain this quantitative conclusion for
the entire branch.

\subsection{Two weak factors give one strong remainder}

Choose a convex neighborhood of the raw state
$(U,\nu,\kappa,a)=(0,0,0,0)$ in the weak product spaces, small in
one finite norm. On it the denominators $D,Q/w$ are positive and
the circle lifts are uniformly increasing. Both Taylor segments
above stay in this neighborhood; the circle map is independent of
$\nu$. All constants below refer to this one neighborhood.

\begin{lemma}[Quadratic remainders in the doubled weight]
\label{lem:posterior-two-weight}
The maps $w^{-1}\mathcal B_{w,U}$ and
$\mathcal B^\circ_{\mathbf v}$ are smooth tame with outputs in,
respectively, the radial space $\mathcal X^0_{2\beta}$ and
$\mathcal K_{2\beta}$. For every fixed number $k$ of state
variations there is a loss $d_k$, independent of the output order
$s$, for which
\begin{equation}\label{eq:posterior-multilinear-bound}
 \begin{split}
 &\|D^k\mathcal B[\eta_1,\ldots,\eta_k]
                         (\eta_{k+1},\eta_{k+2})\|_{2\beta,s}\\
 &\quad\le C_{s,k}\left[
   \sum_{i=1}^{k+2}\|\eta_i\|_{s+d_k}
                    \prod_{j\ne i}\|\eta_j\|_{b+d_k}
   +(1+\|\mathbf y\|_{s+d_k})
                    \prod_j\|\eta_j\|_{b+d_k}\right].
 \end{split}
\end{equation}
The inputs use their weak product norms; compact scalar directions
use their absolute values. The same bound holds for all higher
derivatives of $\mathcal C$ with at least two arguments.
\end{lemma}

\begin{proof}
On component $m$ every control or chart increment, and every positive
height jet, has a factor $H_m^\beta$. Constants at an endpoint cause
no difficulty: their positive jets are supported in the compact
positive region, and vanish on all negative control supports.
Writing $u=H^\beta\widetilde u$ gives, for example,
\[
 H^{-\beta}\mathsf D u
 =\mathsf D\widetilde u+\beta(\mathsf D\log H)\widetilde u.
\]
Thus positive differentiation does not remove a weak factor.
Products and restriction to normalized pieces cost fixed derivatives
and reference factors $\mathcal E_m^{C_{s,k}}$.
Retain two weak factors for the output weight, place the highest
derivative on one input or one state coefficient, and absorb the
\emph{complete} coefficient product into one extra $M$ order.
Finite overlap and
$\|ab\|_{\ell^2}\le\|a\|_{\ell^2}\|b\|_{\ell^\infty}$
sum the estimates. It remains to check that the nonlinear operations
retain these two factors.

For the period this is visible from its integrand. With
$k=w+U_r$, $p=U_\theta/r$ and $E=1+w^2$, it is
\[
 \frac{1+kp_\theta/(1-p^2)}{\sqrt{E(1-p^2)-k^2}}-1.
\]
It is a smooth function of $w$ and the positive jets
$(U_r,U_\theta/r,U_{\theta\theta}/r)$ on the chosen neighborhood.
Each second derivative contains one such jet from each argument.
Further state differentiation cannot remove either argument.
The reciprocal $w^{-1}=w_*^{-1}
 e^{-\nu\chi_*-\sum\kappa_m\chi_m}$ is included in the coefficient
product before its absorption; it does not cost an endpoint exponent.

For the circle map consider a local normalized characteristic flow
$\dot\Theta=G(t,\Theta;\mathbf y)$, using
\eqref{eq:fixed-strip-root-flow} at a root. Its first and second
state variations satisfy
\[
 \dot\Theta_i=G_\theta\Theta_i+G_i,
\]
\[
 \dot\Theta_{12}=G_\theta\Theta_{12}
   +G_{\theta\theta}\Theta_1\Theta_2
   +(G_1)_\theta\Theta_2+(G_2)_\theta\Theta_1+G_{12}.
\]
The homogeneous integrating factor is uniformly bounded, since the
reference speed has zero angular derivative and the perturbation has
a small integrated derivative. $G_i$ has one weak factor and
$G_{12}$ has two. Variation of constants gives one factor in
$\Theta_i$ and two in $\Theta_{12}$. Composing flows does not change
the component weight.

Inversion preserves this count. If $\iota=h^{-1}$, then
\[
 \iota_i=-\frac{h_i\circ\iota}{h'\circ\iota},\qquad
 \iota_{12}=-\frac{h_{12}\circ\iota
 +(h_1'\circ\iota)\iota_2+(h_2'\circ\iota)\iota_1
 +(h''\circ\iota)\iota_1\iota_2}{h'\circ\iota}.
\]
The denominator is bounded below, and every term in the second
numerator contains both variations. Apply this to $f^{-1}$ and
$h_R^{-1}$ in \eqref{eq:full-height-one-chart}; parity projection
and averaging preserve the weights.

For higher variations label the directions by a set $I$. In the
chain rule for $G(t,\Theta;\mathbf y)$ a subset $J\subset I$ acts
directly on the state and a partition of $I\setminus J$ acts on
the flow. Each term is
\[
 (D_{\mathbf y}^{|J|}\partial_\theta^{\ell}G
                   [\eta_j:j\in J])\circ\Theta
                            \prod_{B}\Theta_B,
\]
where $\ell$ is the number of blocks. Keep $G_\theta\Theta_I$
on the linear side; all other terms use proper lower variations and
contain each label exactly once. Induction retains both designated
argument factors. The inverse identity has the same induction.
Additional weak factors are bounded in low norms. Local tame product
and composition estimates put only one factor at high order. The
coefficient formulas, root regularization and finite restrictions
have fixed order; $k$ variations add a cost depending on $k$, not
on $s$. This proves \eqref{eq:posterior-multilinear-bound}.

Finally these are estimates for the actual differentiated formulas.
Taylor's integral remainder with one additional variation proves
continuity and differentiability in the $2\beta$ norm. Iterating
proves smoothness in the stronger Fr\'echet topology, not merely
pointwise stronger decay.
\end{proof}

\subsection{Integrating the mean without losing the exponent}

For $\gamma>0$ a primitive from its own endpoint preserves exponent
$\gamma$. To see this, orient $\varrho$ toward the endpoint. Then
$(\log H)'\to-1$ and $(\log E_0)'\to0$. For a fixed $A$,
the ratio of $H^\gamma E_0^A$ at $\sigma\ge\varrho$ to its value at
$\varrho$ is eventually bounded by
$e^{-\gamma(\sigma-\varrho)/2}$. The conjugated integration kernel
has bounded row and column integrals; $dr=\ell\,d\varrho$ supplies
only a fixed coefficient factor. The $L^2$ Schur estimate applies.
Further derivatives use the integrand itself, while for $c\ge0$
the weight ratio $M(\varrho)^c/M(\sigma)^c\le1$.
The remaining compact interval has ordinary bounded integration.
This proves a tame primitive estimate with fixed loss. For a
primitive normalized at the right endpoint, subtract its left
constant times $\chi_L$ to use the left endpoint estimate.

Let $\mathscr P$ denote this right-normalized primitive. Formula
\eqref{eq:mean-recovered-derivative} identifies the weak mean with
the explicit stronger-space map
\begin{equation}\label{eq:rewrite-mean-recovery}
 \mathcal U(\mathbf x,q)
 =-q\mathscr P[w^{-1}\mathcal B_{w,qu}(u,u)]
 =u^{\rm av}\quad\text{on the solution branch}.
\end{equation}
It is smooth into $\mathcal X_{2\beta}$ and has size $O(q)$ in
every norm. In particular its endpoint scalar is also $O(q)$.
The fixed row \eqref{eq:rewrite-reference-circle-row} maps its
strong radial part into $\mathcal K_{2\beta}$ with fixed loss,
and annihilates the endpoint step. The circle identity therefore gives
\begin{equation}\label{eq:rewrite-control-recovery}
 \kappa_q=-qR_0^\circ\mathcal U(\mathbf x(q),q)
            -\mathcal B^\circ_{\mathbf v(q)}(\mathbf v(q),\mathbf v(q)).
\end{equation}
The weak branch has $\mathbf v(q)=O(q)$. The two terms on the right
are consequently $O(q^2)$ in every $\mathcal K_{2\beta}$ norm.
This identity also proves smoothness of the existing controls in
that stronger topology and $\dot\kappa_0=0$.

The compact scalar is quadratic as well. Differentiate the scalar
equation at $q=0$:
\[
 \mathscr L_{w_*}\dot u_0+\dot\nu_0 k_*+K_*\dot\kappa_0
       +\left.\partial_q\mathcal F_q(w_*,\Phi_*)\right|_{q=0}=0,
 \qquad \ell_*\dot u_0=0.
\]
The final term is quadratic in second-mode jets. Products of
$\cos2\theta$ and $\sin2\theta$ have frequencies zero and four,
so its second-mode projection vanishes. Since $\dot\kappa_0=0$,
the second-mode equation is the homogeneous bordered reference
equation. Its left inverse gives $\dot\nu_0=0$, hence $\nu_q=O(q^2)$.

\begin{figure}[htbp]
\centering
\begin{tikzpicture}[x=6.6cm,y=1cm,font=\small,>=Stealth]
\draw[->] (0,0)--(1.75,0) node[right] {endpoint exponent};
\fill[blue!8] (.5,-.12) rectangle (1,.12);
\foreach \x/\lab in {0/0,.5/{1/2},.75/{\beta=3/4},1/1,1.5/{2\beta=3/2}}
 {\draw (\x,-.12)--(\x,.12); \node[below] at (\x,-.16) {$\lab$};}
\draw[->,thick,orange!85!black] (.75,.4) to[bend left=30] (1.5,.4);
\node[above] at (1.12,.95) {two weak factors in each mean remainder};
\node[align=center] at (.75,-1.05) {choice for both estimates\\$1/2<\beta<1$};
\end{tikzpicture}
\caption{The choice $\beta=3/4$ serves two estimates. Neighboring
exponents below one control the reference inverse. The mean equations
multiply two weak factors and produce an exponent above one, giving
the strengthened metric bound in Theorem~\ref{thm:new-annular-height}.
Flat extension already follows from any positive exponent.}
\label{fig:exponent-gain}
\end{figure}

\subsection{The annular conclusion}

\begin{theorem}[A nonradial height with flat ends]
\label{thm:new-annular-height}
For the fixed reference of Proposition~\ref{prop:rad-fixed-theorem}
there is $q_0>0$ and a smooth branch
$(u_q,\nu_q,\kappa_q,h_{L,q})$, $|q|<q_0$, through
$(\Phi_*,0,0,\id)$, with the following properties.
Set $w_q=w_*e^{\nu_q\chi_*+\sum_m\kappa_{q,m}\chi_m}$ and
$z_q=h_{w_q}+qu_q$, choosing a fixed additive normalization for
$h_{w_q}'=w_q$. Then
\[
 \mathcal N_{h_{w_q}}(z_q)=0,\qquad
 |dz_q|_{(1+w_q^2)dr^2+r^2d\theta^2}<1,\qquad \ell_*u_q=1.
\]
The profile is positive in the annular interior and extends smoothly
by zero across both endpoints. The height is even and $\pi$-periodic
in angle, extends smoothly by a constant at each endpoint, and is
nonradial for $q\ne0$. The two constants may differ.
All roots of $w_q'$ and their radial germs agree with the reference.
The two characteristic returns of the true height derivative obey
\begin{equation}\label{eq:annular-return-condition}
 h_{R,q,m}f_{q,m,\pm}h_{L,q,m}^{-1}=R_{\pm T_m}
\end{equation}
for increasing odd $\pi$-equivariant charts and the fixed reference angles.
On the same interval of $q$, for every $s$,
\begin{equation}\label{eq:annular-branch-estimates}
 \|u_q-\Phi_*\|_{\mathcal X_\beta,s}\le C_s|q|,\qquad
 |\nu_q|+\|\kappa_q\|_{\mathcal K_{2\beta},s}\le C_sq^2.
\end{equation}
The branch is smooth with these topologies. In particular, for
every $a\ge0$ and $\epsilon>0$,
\[
 |\partial_r^a(w_q-w_*)|\le C_{a,\epsilon}q^2H^{1-\epsilon}
\]
at the endpoints, and $w_q-w_*=O(q^2)$ in every smooth norm.
\end{theorem}

\begin{proof}
The weak branch was constructed in Section~\ref{sec:solving-the-equation}.
Equations \eqref{eq:rewrite-mean-recovery} and
\eqref{eq:rewrite-control-recovery} give the strengthened estimates
and smooth dependence just proved. Choose $q_0$ so that
$[-q_0,q_0]$ lies inside its domain. Every seminorm of this smooth
branch is bounded on that one compact interval; the constants depend
on the order, but $q_0$ does not.

On the tails the compact control vanishes. Differentiating the
formula for $w_q$ gives
$|\partial_r^a(w_q-w_*)|\le C_aq^2H^{2\beta}E_0^{C_a}$.
Since $2\beta>1$, \eqref{eq:full-material-envelope} implies the
displayed endpoint estimate. In the compact middle use smoothness
and the quadratic bound for $\nu_q$. All controls vanish near each
turning circle, preserving its full metric germ; the common low
neighborhood preserves the signs elsewhere.

At a zero both residual corrections vanish. The exact return
identities for $\widehat L$ therefore hold for the true derivative.
The common inequality $D>0$ is equivalent to $|dz_q|_g<1$.
The scalar residual vanishes, and the curvature formula identifies
this with flatness of $g-dz_q^2$. Since $\ell_*u_q=1$, the second
angular component is nonzero, so $z_q$ is nonradial for $q\ne0$.

The height space already gives a flat extension after subtracting
its endpoint constants, including all parameter derivatives.
Extend $w_q$ by zero and $z_q$ by those constants and enlarge the
annulus slightly. The enlarged annulus has zero-profile open collars around its
original positive-profile core. The remainder metric stays positive
and flat on these collars, giving the data needed for geometric
closure. This proves the theorem.
\end{proof}

\subsection{Fixing the metric and moving the surface}

We finish by distinguishing the two parameters. The small parameter
$q$ was used to construct one nonradial realization of one rotational
metric. Choose a sufficiently small nonzero value and keep it fixed.
The annular theorem then supplies exactly the data of
Theorem~\ref{thm:geometric-closure}: a positive flat metric
$g-dz_q^2$, the half-turn symmetry, and two constant-height collars.

That theorem develops the graph, fills its inner disk, and attaches
the fixed marked complementary cap. Its independent phase parameter
$\tau$ gives the jointly smooth periodic family
\[
 F_\tau(r,\theta)=R_{-2\pi\tau}F(r,\theta+2\pi\tau)
\]
on the moving annulus and the fixed map elsewhere. The collars cancel
the two rotations pointwise. The marked cap makes two resulting
closed images noncongruent. The same geometric theorem proves
embeddedness, constant area, constant enclosed volume and constant
total oriented mean curvature. This proves Theorem~\ref{thm:main}.
